% 褐矮星（综述）
% license CCBYSA3
% type Wiki

本文根据 CC-BY-SA 协议转载翻译自维基百科\href{https://en.wikipedia.org/wiki/Brown_dwarf}{相关文章}。

\begin{figure}[ht]
\centering
\includegraphics[width=6cm]{./figures/6432b0e71e12ceea.png}
\caption{T 型褐矮星的艺术概念图} \label{fig_Brdw_1}
\end{figure}
\begin{figure}[ht]
\centering
\includegraphics[width=6cm]{./figures/dc30c14f9295e490.png}
\caption{对比：大多数褐矮星的体积仅比木星略大（约大15–20\%）$^{[1]}$，但由于其密度更高，质量仍可达到木星的 80 倍。图像按比例绘制，其中木星的半径约为地球的 11 倍，而太阳的半径约为木星的 10 倍。} \label{fig_Brdw_2}
\end{figure}
褐矮星是亚恒星天体，其质量大于最大的气态巨行星，但小于质量最小的主序星。它们的质量大约为木星的 13–80 倍（($M_J$)）$^{[2]}$——不足以在其核心中维持将氢聚变为氦的核反应，但又足够巨大，可以通过氘（(^2\mathrm{H})）的聚变释放一定的光和热。氘是氢的一种同位素，含有一个中子和一个质子，其聚变所需温度更低。质量最大的褐矮星（$(>65,M_J)$）还能够进行锂（$(^7\mathrm{Li})$）的聚变。

天文学家按照光谱型对自发光天体进行分类，这一分类与表面温度密切相关。褐矮星分布在 M 型（2100–3500 K）、L 型（1300–2100 K）、T 型（600–1300 K）和 Y 型（<600 K） 光谱类别中 $^{[3][4]}$。由于褐矮星不会发生稳定的氢聚变，它们会随着时间推移逐渐冷却，并在演化过程中依次进入更晚的光谱型。

“褐矮星”中的“褐色”原意是指介于红色与黑色之间的一种颜色 $^{[5]}$。以肉眼观察，大多数褐矮星呈现为品红色，而其他褐矮星则会因温度不同而显示出不同的颜色 $^{[3][6]}$。褐矮星可能是完全对流的，内部不存在明显的分层结构或随深度变化的化学分异 $^{[7]}$。

尽管早在 20 世纪 60 年代就有人从理论上预言了褐矮星的存在，但直到 1994 年，第一批明确无误的褐矮星才被发现 $^{[8]}$。由于褐矮星的表面温度相对较低，它们在可见光波段并不明亮，而主要在红外波段辐射能量。随着红外探测设备能力的提升，天文学家已经识别出数以千计的褐矮星。目前已知距离最近的褐矮星位于 Luhman 16 系统，这是一个由 L 型和 T 型褐矮星组成的双星系统，距离太阳约 6.5 光年（2.0 秒差距）。在距离上，Luhman 16 是继半人马座 $\alpha$ 星系和巴纳德星之后，距离太阳第三近的恒星系统。
\subsection{历史}
